% === Begin Label Data ===
% ID: 196
% 难度星级: 
% 标签: 
% 备注: 
% 组卷引用次数: 0
% === End  Label Data ===

\begin{problem}{2012}{G}{上海卷（文）}{23}{数列}
对于项数为 $m$ 的有穷数列 $\{a_n\}$, 记 $b_k = \max \{a_1, a_2, \cdots, a_k\}$ ($k=1, 2, \cdots, m$), 即 $b_k$ 为 $a_1, a_2, \cdots, a_k$ 中的最大值, 并称数列 $\{b_n\}$ 是 $\{a_n\}$ 的控制数列, 如 1, 3, 2, 5, 5 的控制数列是 1, 3, 3, 5, 5.

（1） 若各项均为正整数的数列 $\{a_n\}$ 的控制数列为 2, 3, 4, 5, 5, 写出所有的 $\{a_n\}$

（2） 设 $\{b_n\}$ 是 $\{a_n\}$ 的控制数列, 满足 $a_k + b_{m-k+1} = C$ ($C$ 为常数, $k=1, 2, \cdots, m$), 求证: $b_k = a_k$ ($k=1, 2, \cdots, m$)

（3） 设 $m=100$, 常数 $a \in \left(\dfrac{1}{2}, 1\right)$, 若 $a_n = a n^2 - (-1)^{\frac{n(n+1)}{2}} n$, $\{b_n\}$ 是 $\{a_n\}$ 的控制数列, 求 $(b_1 - a_1) + (b_2 - a_2) + \cdots + (b_{100} - a_{100})$
\end{problem}

\begin{solutions}
（1） 数列 $\{a_n\}$ 为: 2, 3, 4, 5, 1;\quad 2, 3, 4, 5, 2;\quad 2, 3, 4, 5, 3;\quad 2, 3, 4, 5, 4;\quad 

2, 3, 4, 5, 5. \hfill $\cdots\cdots 4$ 分

（2） 因为 $b_k = \max \{a_1, a_2, \cdots, a_k\}$, $b_{k+1} = \max \{a_1, a_2, \cdots, a_k, a_{k+1}\}$.

所以 $b_{k+1} \ge b_k$. \hfill $\cdots\cdots 6$ 分

因为 $a_k + b_{m-k+1} = C$, $a_{k+1} + b_{m-k} = C$.

所以 $a_{k+1} - a_k = b_{m-k+1} - b_{m-k} \ge 0$, 即 $a_{k+1} \ge a_k$. \hfill $\cdots\cdots 8$ 分

因此, $b_k = a_k$. \hfill $\cdots\cdots 10$ 分

（3） 对 $k=1, 2, \cdots, 25$,

$a_{4k-3} = a(4k-3)^2 + (4k-3)$;\quad   $a_{4k-2} = a(4k-2)^2 + (4k-2)$;

$a_{4k-1} = a(4k-1)^2 - (4k-1)$;\quad   $a_{4k} = a(4k)^2 - (4k)$. \hfill $\cdots\cdots 12$ 分

比较大小, 可得 $a_{4k-2} > a_{4k-3}$.

因为 $\dfrac{1}{2} < a < 1$, 所以 $a_{4k-1} - a_{4k-2} = (a-1)(8k-3) < 0$, 即 $a_{4k-2} > a_{4k-1}$;

$a_{4k} - a_{4k-2} = 2(2a-1)(4k-1) > 0$, 即 $a_{4k} > a_{4k-2}$.

又 $a_{4k+1} > a_{4k}$.

从而 $b_{4k-3} = a_{4k-3}$, $b_{4k-2} = a_{4k-2}$, $b_{4k-1} = a_{4k-2}$, $b_{4k} = a_{4k}$. \hfill $\cdots\cdots 15$ 分

因此 
\begin{align*}
& (b_1 - a_1) + (b_2 - a_2) + \cdots + (b_{100} - a_{100})  \\
=& (a_2 - a_3) + (a_6 - a_7) + \cdots + (a_{98} - a_{99})   \\ 
=& \sum_{k=1}^{25} (a_{4k-2} - a_{4k-1})\\
=& (1-a) \sum_{k=1}^{25} (8k-3) = 2525(1-a).
\end{align*}
 \hfill $\cdots\cdots 18$ 分
\end{solutions}
